\chapter{Dados por item do \textit{bootstrap} no conjunto de teste}
\label{ap:bootstrap}

Este apêndice apresenta os dados de que depende o \textit{bootstrap} do
conjunto de teste (Seções~\ref{subsec:etapaD} e~\ref{sec:res_teste}). Para cada
um dos 26 itens rotulados, registraram-se quatro indicadores: se o item é um
edital descartado como listagem ($p_i$, edital perdido) e se foi aceito sem
ser um edital aberto e relevante ($e_i$, item errado aceito), na configuração
original ($o$) e na calibrada ($c$). Os indicadores da configuração original
foram obtidos com o código da versão anterior à calibração, pois parte das
regras daquela configuração foi removida do código.

A Tabela~\ref{tab:bootstrap_itens} lista os 9 itens com pelo menos um
indicador igual a 1. Nos 17 itens restantes, os quatro indicadores são nulos,
e esses itens só contribuem para o \textit{bootstrap} com diferença zero. Os
indicadores dos 26 itens estão no arquivo \texttt{bootstrap\_itens.csv} do
diretório da coleta de teste, o que permite refazer o cálculo da
Equação~\ref{eq:bootstrap} sem o código do sistema.

\begin{table}[htbp]
\centering
\caption{Itens do conjunto de teste com algum indicador não nulo}
\label{tab:bootstrap_itens}
\footnotesize
\setlength{\tabcolsep}{4pt}
\begin{tabularx}{\textwidth}{l >{\raggedright\arraybackslash}X c c c c c}
\toprule
 & & & \multicolumn{2}{c}{Perdido} & \multicolumn{2}{c}{Errado} \\
\cmidrule(lr){4-5} \cmidrule(lr){6-7}
Item & Título & Edital & $p_i^{o}$ & $p_i^{c}$ & $e_i^{o}$ & $e_i^{c}$ \\
\midrule
\texttt{74b6d263} & Edital de seleção MICSUL 2026 --- Ministério da Cultura & sim & 1 & 0 & 0 & 0 \\
\texttt{7e538218} & Chamada Pública 13/2026 - Ipea & sim & 1 & 0 & 0 & 0 \\
\texttt{f2aa1063} & Chamada de Propostas 42/2026: Auxílio à Inovação Regular\ldots{} & sim & 0 & 0 & 1 & 0 \\
\texttt{aba75096} & Chamada Pública 008/2026 - Ipea & sim & 1 & 0 & 0 & 0 \\
\texttt{10dcced2} & Chamada Pública nº 027/2026 - Ipea & sim & 1 & 0 & 0 & 0 \\
\texttt{59f6fdcd} & Chamada de Propostas FAPESP e Ohio State University 2026 & sim & 0 & 0 & 0 & 1 \\
\texttt{c49162cf} & [PDF] licitacao@cnpq.br . CONDIÇÕES PARA RETIRADA DO EDIT\ldots{} & sim & 1 & 1 & 0 & 0 \\
\texttt{0e026f2a} & Edital MCT / CNPq No 70/2008 & sim & 0 & 0 & 1 & 0 \\
\texttt{63adc331} & Prêmio Jovem Cientista | FRM - Fundação Roberto Marinho & não & 0 & 0 & 0 & 1 \\
\midrule
\multicolumn{3}{l}{Total nos 26 itens} & 5 & 1 & 2 & 2 \\
\bottomrule
\end{tabularx}
\fonte{Elaborado pelo autor. Os identificadores estão abreviados aos oito
primeiros caracteres, e os títulos longos, truncados.}
\end{table}

A tabela mostra a origem das duas diferenças. Nos editais perdidos, a
configuração calibrada recuperou quatro editais e não perdeu nenhum que a
original não perdesse; as diferenças por item, $p_i^{c} - p_i^{o}$, são todas
nulas ou negativas, e por isso nenhuma reamostra pode indicar piora. O único
edital perdido pelas duas configurações é o item \texttt{c49162cf}, rotulado
como aberto sem data-limite explícita. Nos itens errados aceitos, o total igual
esconde uma troca: as duas configurações aceitaram indevidamente itens
diferentes. A original aceitou dois editais com inscrições encerradas, ambos
fora do tema, e a calibrada, um edital com inscrições encerradas e a página de
listagem de um prêmio. As diferenças por item são $+1$ em dois itens e $-1$ em outros dois,
e o resultado de cada reamostra depende de quantos de cada grupo são
sorteados, o que explica o intervalo largo e centrado em zero.
